% GENERATED FILE. Edit canonical lessons, then run npm run generate:books.
% canonical-source-hash: fnv1a64:34d77bc321fdd583
% canonical-lessons: TE-C409-himsa, TE-C409-samarasyam, TE-C409-aikyata, TE-C409-sanghibhavam, TE-C409-acarincu

\chapter{Living Together}
\label{ch:te-living-together}
\hlchaptermodality{409}

\begin{chapteropening}
\textbf{By the end of this chapter:} \emph{I can talk about violence, harmony, unity and solidarity.}

Complete the last lesson of chapter 409: I can talk about violence, harmony, unity and solidarity.
\end{chapteropening}

\section[\texorpdfstring{\te{హింస} (hiṁsa)}{హింస (hiṁsa)}]{\te{హింస} (hiṁsa) — violence}

\label{lesson:TE-C409-himsa}



\begin{quote}
Before the new one: say the Telugu for a member, then the Telugu for discrimination.
\end{quote}

\subsection*{You'll want to know: \te{హింస}}
\textbf{\te{హింస}} — \emph{hiṁsa} — \textquotedblleft{}violence\textquotedblright{}.

\begin{quote}
\te{హింస} \te{ఏ} \te{సమస్యకూ} \te{పరిష్కారం} \te{కాదు}.
\end{quote}

\begin{quote}
\emph{hiṁsa ē samasyakū pariṣkāraṁ kādu.}
\end{quote}

\begin{quote}
Violence is no solution to any problem.
\end{quote}

\te{అహింస}, non-violence.

\begin{cousinweb}
\textbf{\te{హింస}} (\emph{hiṁsa}) is borrowed from Sanskrit \emph{hiṁsā}, \textquotedblleft{}injury, harm, hurt\textquotedblright{} (Monier-Williams).
\end{cousinweb}

\begin{grammarlens}[title={What you've built}]
One more word for giving reasons and opinions.
\end{grammarlens}

\subsection*{Your turn}
\begin{itemize}
\raggedright
  \item \emph{Say it:} \emph{hiṁsa}
  \item \emph{Say it:} \emph{hiṁsa}, once more
  \item \emph{Say it:} \emph{hiṁsa ē samasyakū pariṣkāraṁ kādu.} — Violence is no solution to any problem.
  \item \emph{Recall:} say the Telugu for a member, then the Telugu for discrimination, then say \emph{hiṁsa} again
\end{itemize}

\begin{tcolorbox}[breakable,colback=teal!4,colframe=teal!35!black,title={Before you move on}]
What does \te{హింస} mean? (\textquotedblleft{}Violence\textquotedblright{}.) Say it once more.
\end{tcolorbox}
\section[\texorpdfstring{\te{సామరస్యం} (sāmarasyaṁ)}{సామరస్యం (sāmarasyaṁ)}]{\te{సామరస్యం} (sāmarasyaṁ) — harmony; concord}

\label{lesson:TE-C409-samarasyam}



\begin{quote}
Before the new one: say the Telugu for discrimination, then the Telugu for violence.
\end{quote}

\subsection*{You'll want to know: \te{సామరస్యం}}
\textbf{\te{సామరస్యం}} — \emph{sāmarasyaṁ} — \textquotedblleft{}harmony; concord\textquotedblright{}.

\begin{quote}
\te{మా} \te{వీధిలో} \te{అందరూ} \te{సామరస్యంగా} \te{ఉంటారు}.
\end{quote}

\begin{quote}
\emph{mā vīdhilō andarū sāmarasyaṅgā uṇṭāru.}
\end{quote}

\begin{quote}
Everyone in our street lives in harmony.
\end{quote}

\te{సామరస్యంగా}, in harmony, peaceably.

\begin{cousinweb}
\textbf{\te{సామరస్యం}} (\emph{sāmarasyaṁ}) is built on Sanskrit \emph{samarasa}, \textquotedblleft{}having equal feelings\textquotedblright{} (Monier-Williams), from \emph{sama}, \textquotedblleft{}same\textquotedblright{}, and \emph{rasa}, \textquotedblleft{}feeling\textquotedblright{}.
\end{cousinweb}

\begin{grammarlens}[title={What you've built}]
One more word for giving reasons and opinions.
\end{grammarlens}

\subsection*{Your turn}
\begin{itemize}
\raggedright
  \item \emph{Say it:} \emph{sāmarasyaṁ}
  \item \emph{Say it:} \emph{sāmarasyaṁ}, once more
  \item \emph{Say it:} \emph{mā vīdhilō andarū sāmarasyaṅgā uṇṭāru.} — Everyone in our street lives in harmony.
  \item \emph{Recall:} say the Telugu for discrimination, then the Telugu for violence, then say \emph{sāmarasyaṁ} again
\end{itemize}

\begin{tcolorbox}[breakable,colback=teal!4,colframe=teal!35!black,title={Before you move on}]
What does \te{సామరస్యం} mean? (\textquotedblleft{}Harmony; concord\textquotedblright{}.) Say it once more.
\end{tcolorbox}
\section[\texorpdfstring{\te{ఐక్యత} (aikyata)}{ఐక్యత (aikyata)}]{\te{ఐక్యత} (aikyata) — unity}

\label{lesson:TE-C409-aikyata}



\begin{quote}
Before the new one: say the Telugu for violence, then the Telugu for harmony.
\end{quote}

\subsection*{You'll want to know: \te{ఐక్యత}}
\textbf{\te{ఐక్యత}} — \emph{aikyata} — \textquotedblleft{}unity\textquotedblright{}.

\begin{quote}
\te{ఐక్యత} \te{లేకపోతే} \te{మనం} \te{గెలవలేం}.
\end{quote}

\begin{quote}
\emph{aikyata lēkapōtē manaṁ gelavalēṁ.}
\end{quote}

\begin{quote}
Without unity we cannot win.
\end{quote}

\te{ఐక్యంగా}, united, together.

\begin{cousinweb}
\textbf{\te{ఐక్యత}} (\emph{aikyata}) is built on Sanskrit \emph{aikya}, \textquotedblleft{}oneness\textquotedblright{} (CDIAL 2533), from \emph{eka}, \textquotedblleft{}one\textquotedblright{}, with the abstract ending -\te{త}.
\end{cousinweb}

\begin{grammarlens}[title={What you've built}]
One more word for giving reasons and opinions.
\end{grammarlens}

\subsection*{Your turn}
\begin{itemize}
\raggedright
  \item \emph{Say it:} \emph{aikyata}
  \item \emph{Say it:} \emph{aikyata}, once more
  \item \emph{Say it:} \emph{aikyata lēkapōtē manaṁ gelavalēṁ.} — Without unity we cannot win.
  \item \emph{Recall:} say the Telugu for violence, then the Telugu for harmony, then say \emph{aikyata} again
\end{itemize}

\begin{tcolorbox}[breakable,colback=teal!4,colframe=teal!35!black,title={Before you move on}]
What does \te{ఐక్యత} mean? (\textquotedblleft{}Unity\textquotedblright{}.) Say it once more.
\end{tcolorbox}
\section[\texorpdfstring{\te{సంఘీభావం} (saṅghībhāvaṁ)}{సంఘీభావం (saṅghībhāvaṁ)}]{\te{సంఘీభావం} (saṅghībhāvaṁ) — solidarity}

\label{lesson:TE-C409-sanghibhavam}



\begin{quote}
Before the new one: say the Telugu for harmony, then the Telugu for unity.
\end{quote}

\subsection*{You'll want to know: \te{సంఘీభావం}}
\textbf{\te{సంఘీభావం}} — \emph{saṅghībhāvaṁ} — \textquotedblleft{}solidarity\textquotedblright{}.

\begin{quote}
\te{రైతులకు} \te{విద్యార్థులు} \te{సంఘీభావం} \te{తెలిపారు}.
\end{quote}

\begin{quote}
\emph{raitulaku vidyārthulu saṅghībhāvaṁ telipāru.}
\end{quote}

\begin{quote}
The students expressed solidarity with the farmers.
\end{quote}

\te{సంఘీభావం} \te{తెలుపు}, to show solidarity; the group takes -\te{కి}.

\begin{cousinweb}
\textbf{\te{సంఘీభావం}} (\emph{saṅghībhāvaṁ}) is built from Sanskrit \emph{saṁgha}, \textquotedblleft{}an association\textquotedblright{} (CDIAL 12854), and \emph{bhāva}, \textquotedblleft{}being, state\textquotedblright{} (CDIAL 9475): the state of standing together.
\end{cousinweb}

\begin{grammarlens}[title={What you've built}]
One more word for giving reasons and opinions.
\end{grammarlens}

\subsection*{Your turn}
\begin{itemize}
\raggedright
  \item \emph{Say it:} \emph{saṅghībhāvaṁ}
  \item \emph{Say it:} \emph{saṅghībhāvaṁ}, once more
  \item \emph{Say it:} \emph{raitulaku vidyārthulu saṅghībhāvaṁ telipāru.} — The students expressed solidarity with the farmers.
  \item \emph{Recall:} say the Telugu for harmony, then the Telugu for unity, then say \emph{saṅghībhāvaṁ} again
\end{itemize}

\begin{tcolorbox}[breakable,colback=teal!4,colframe=teal!35!black,title={Before you move on}]
What does \te{సంఘీభావం} mean? (\textquotedblleft{}Solidarity\textquotedblright{}.) Say it once more.
\end{tcolorbox}
\section[\texorpdfstring{\te{ఆచరించు} (ācariñcu)}{ఆచరించు (ācariñcu)}]{\te{ఆచరించు} (ācariñcu) — to put into practice; to observe}

\label{lesson:TE-C409-acarincu}



\begin{quote}
Before the new one: say the Telugu for unity, then the Telugu for solidarity.
\end{quote}

\subsection*{You'll want to know: \te{ఆచరించు}}
\textbf{\te{ఆచరించు}} — \emph{ācariñcu} — \textquotedblleft{}to put into practice; to observe\textquotedblright{}.

\begin{quote}
\te{అతను} \te{చెప్పిన} \te{మాటను} \te{తానే} \te{ఆచరించాడు}.
\end{quote}

\begin{quote}
\emph{atanu ceppina māṭanu tānē ācariñcāḍu.}
\end{quote}

\begin{quote}
He practised what he preached.
\end{quote}

Past (I …): \te{ఆచరించాను}, \emph{ācariñcānu}.

For customs, rules and advice: \te{ఆచారం}, a custom, which you know, is the same family.

\begin{cousinweb}
\textbf{\te{ఆచరించు}} (\emph{ācariñcu}) is built on Sanskrit \emph{ācaraṇa}, \textquotedblleft{}undertaking, practising, performing; conduct\textquotedblright{} (Monier-Williams), with the Telugu verb ending -\te{ఇంచు}.
\end{cousinweb}

\begin{grammarlens}[title={What you've built}]
One more word for giving reasons and opinions.
\end{grammarlens}

\subsection*{Your turn}
\begin{itemize}
\raggedright
  \item \emph{Say it:} \emph{ācariñcu}
  \item \emph{Say it:} \emph{ācariñcu}, once more
  \item \emph{Say it:} \emph{atanu ceppina māṭanu tānē ācariñcāḍu.} — He practised what he preached.
  \item \emph{Recall:} say the Telugu for unity, then the Telugu for solidarity, then say \emph{ācariñcu} again
\end{itemize}

\begin{tcolorbox}[breakable,colback=teal!4,colframe=teal!35!black,title={Before you move on}]
What does \te{ఆచరించు} mean? (\textquotedblleft{}To put into practice; to observe\textquotedblright{}.) Say it once more.
\end{tcolorbox}
